\chapter{DEFINITIONS}
\label{app:definitions}

\section*{A.1 Network Architecture Definitions}
\addcontentsline{toc}{section}{A.1 Network Architecture Definitions}

\subsection*{Spine Switch}
\label{def:spine}
A \textbf{spine switch} forms the core forwarding layer in spine-leaf architecture. 
It connects only to leaf switches (no direct server connections) and operates 
at Layer 3 for routing between leaf switches. In our design, EoR-1 and EoR-2 
serve as spine switches.

\subsection*{Leaf Switch (ToR)}
\label{def:leaf}
A \textbf{leaf switch} (Top of Rack) provides server connectivity in the access 
layer. It connects to servers via access ports and uplinks to all spine switches 
via trunk ports. In our design, we have eight leaf switches (ToR A-1 through 
ToR D-2).

\subsection*{Spine-Leaf Architecture}
\label{def:spineleaf}
\textbf{Spine-Leaf architecture} is a two-tier network topology where every 
leaf switch connects to every spine switch in a full mesh. Any server reaches 
any other server in exactly two hops, providing consistent latency and ECMP 
load balancing.

\subsection*{Three-Tier Architecture}
\label{def:threetier}
\textbf{Three-Tier architecture} is a traditional data center architecture 
consisting of core layer (backbone), aggregation layer (policy enforcement), 
and access layer (server connectivity). Designed for north-south traffic.

\subsection*{EoR (End of Row)}
\label{def:eor}
An \textbf{End-of-Row switch} is placed at the end of a row of server racks, 
connecting all servers in that row. In our design, EoR-1 and EoR-2 function as 
spine switches.

\subsection*{ToR (Top of Rack)}
\label{def:tor}
A \textbf{Top-of-Rack switch} is located at the top of a server cabinet, 
connecting all servers within that cabinet. In our design, we have eight 
ToR switches (two per cabinet).

\subsection*{Full Mesh}
\label{def:fullmesh}
\textbf{Full mesh} is a topology where every device connects directly to every 
other device. In spine-leaf architecture, every leaf switch connects to every 
spine switch.
\subsection*{Cluster}
\label{def:cluster}
A \textbf{cluster} is a group of physical servers that act as a single 
logical unit to provide high availability and load balancing. In our 
project, one Linux instance represents a cluster of four physical servers.

\section*{A.2 VLAN and Switching Definitions}
\addcontentsline{toc}{section}{A.2 VLAN and Switching Definitions}

\subsection*{VLAN (Virtual Local Area Network)}
\label{def:vlan}
A \textbf{VLAN} is a logical subdivision of a physical network that groups 
devices based on function or security requirements rather than physical 
location. VLANs isolate traffic, reduce broadcast domains, and improve 
security.

\subsection*{Trunk Port}
\label{def:trunk}
A \textbf{trunk port} is a switch port that carries traffic for multiple 
VLANs simultaneously using 802.1Q tagging. Used for inter-switch connections.

\subsection*{Access Port}
\label{def:accessport}
An \textbf{access port} is a switch port that belongs to a single VLAN and 
carries untagged traffic. Used for connecting end devices like servers.

\section*{A.3 Redundancy and High Availability Definitions}
\addcontentsline{toc}{section}{A.3 Redundancy and High Availability Definitions}

\subsection*{VRRP (Virtual Router Redundancy Protocol)}
\label{def:vrrp}
\textbf{VRRP} is a protocol that allows multiple routers to share a single 
virtual IP address. If the master router fails, a backup takes over 
automatically within seconds, providing gateway redundancy.

\subsection*{MLAG (Multi-Chassis Link Aggregation)}
\label{def:mlag}
\textbf{MLAG} is a technology that allows two physical switches to appear as 
a single logical switch to connected devices. In our design, MLAG pairs operate 
in active-standby mode with master/backup switches.

\subsection*{Active-Active}
\label{def:activeactive}
\textbf{Active-active} is an alternative redundancy mode where both devices forward traffic 
simultaneously. This mode is not used in our design.

\subsection*{Active-Standby}
\label{def:activestandby}
\textbf{Active-standby} is a redundancy mode where one device (master) handles 
all traffic while the other (standby) remains idle until the master fails. 
VRRP uses this mode for gateways and firewalls. MLAG pairs at the ToR and EoR 
layers also operate in active-standby mode in our design.

\subsection*{Failover}
\label{def:failover}
\textbf{Failover} is the automatic switching from a failed primary component 
to a redundant backup component.

\section*{A.4 Security Definitions}
\addcontentsline{toc}{section}{A.4 Security Definitions}

\subsection*{DMZ (Demilitarized Zone)}
\label{def:dmz}
\textbf{DMZ} is a semi-trusted zone for public-facing services such as web 
servers. It sits between the internet (untrust) and the internal network 
(trust). Accessible from internet on specific ports (HTTP/HTTPS).

\subsection*{Trust Zone}
\label{def:trust}
The \textbf{Trust zone} is the most secure zone for internal services including 
application servers, database servers, and management networks. No direct 
internet access.

\subsection*{Untrust Zone}
\label{def:untrust}
The \textbf{Untrust zone} represents all external networks (internet) and is 
considered completely untrusted. All traffic must pass through firewall 
inspection.

\subsection*{Firewall}
\label{def:firewall}
A \textbf{firewall} is a security device that inspects and filters network 
traffic based on configured rules. Enforces zone-based security policies and 
performs stateful packet inspection.

\subsection*{NAT (Network Address Translation)}
\label{def:nat}
\textbf{NAT} is a method that maps private IP addresses to public IP addresses, 
allowing internal devices to communicate with the internet while hiding 
internal network structure.

\section*{A.5 Virtualization Definitions}
\addcontentsline{toc}{section}{A.5 Virtualization Definitions}

\subsection*{Hypervisor}
\label{def:hypervisor}
A \textbf{hypervisor} is software that creates and runs virtual machines by 
abstracting the underlying physical hardware (e.g., VMware ESXi).

\subsection*{Virtual Machine (VM)}
\label{def:vm}
A \textbf{virtual machine} is an emulation of a computer system that runs its 
own operating system on top of a hypervisor.

\subsection*{Server Cluster}
\label{def:servercluster}
A \textbf{server cluster} is a group of physical servers that act as a single 
logical unit for high availability and load balancing. In our project, one 
Linux instance represents a cluster.

\subsection*{VMware}
\label{def:vmware}
\textbf{VMware} is a virtualization platform that allows running multiple 
virtual machines on a single physical server.

\section*{A.6 Simulation and Tools Definitions}
\addcontentsline{toc}{section}{A.6 Simulation and Tools Definitions}

\subsection*{eNSP (Enterprise Network Simulation Platform)}
\label{def:ensp}
\textbf{eNSP} is Huawei's official network simulation platform for designing 
and testing networks without physical hardware. Requires VirtualBox 5.2.44.

\subsection*{PNETLab (Practical Networking Lab)}
\label{def:pnetlab}
\textbf{PNETLab} is a community-maintained network simulation platform 
(EVE-NG fork) that runs on VMware and supports multiple device images 
(Arista vEOS, Debian Linux, Huawei CE12800).

\subsection*{Arista vEOS}
\label{def:veos}
\textbf{Arista vEOS} (Virtual Extensible Operating System) is Arista's virtual 
network device emulator used in our project for spine and leaf switches. 
Supports VLANs, trunk ports, VRRP, OSPF, and MLAG with 2GB RAM per instance.

\subsection*{VMware Workstation Pro}
\label{def:vmwareworkstation}
\textbf{VMware Workstation Pro} is a Type-2 hypervisor used to host the 
PNETLab virtual machine, providing hardware virtualization support.

\subsection*{PuTTY}
\label{def:putty}
\textbf{PuTTY} is a free SSH and Telnet client used in our project to access 
network devices and servers for configuration and testing.

\subsection*{VPCS (Virtual PC Simulator)}
\label{def:vpcs}
\textbf{VPCS} is a lightweight tool included with PNETLab that simulates a 
simple PC. Limited to one IP address per device, replaced by Linux Tiny Core.

\section*{A.7 Server and Network Interface Definitions}
\addcontentsline{toc}{section}{A.7 Server and Network Interface Definitions}

\subsection*{NIC (Network Interface Card)}
\label{def:nic}
A \textbf{NIC} (Network Interface Card) is a hardware component that connects 
a server to a network. In our design, each server has four NICs (M1, M2, S1, S2).

\subsection*{Management Network (M1/M2)}
\label{def:management}
The \textbf{management network} is the network interface used by administrators 
to access servers for configuration, monitoring, and maintenance. Isolated 
from service traffic.

\subsection*{Service Network (S1/S2)}
\label{def:service}
The \textbf{service network} is the network interface used for application 
traffic between servers (web to app, app to database). Separate from 
management network for security.

\section*{A.8 Documentation Definitions}
\addcontentsline{toc}{section}{A.8 Documentation Definitions}

\subsection*{HLD (High-Level Design)}
\label{def:hld}
\textbf{HLD} (High-Level Design) is a top-level design document showing the 
overall architecture, major components, and their relationships without 
implementation details.

\subsection*{LLD (Low-Level Design)}
\label{def:lld}
\textbf{LLD} (Low-Level Design) is a detailed design document containing 
specific IP addresses, VLAN assignments, port connections, configuration 
commands, and implementation procedures.

\subsection*{UML (Unified Modeling Language)}
\label{def:uml}
\textbf{UML} (Unified Modeling Language) is a standardized modeling language 
used to visualize system design. In our project, we used use case diagrams 
and sequence diagrams.

\subsection*{Use Case Diagram}
\label{def:usecase}
A \textbf{use case diagram} is a UML diagram showing actors (users or systems) 
and their interactions with the system. Used to identify functional 
requirements.

\subsection*{Sequence Diagram}
\label{def:sequence}
A \textbf{sequence diagram} is a UML diagram showing the temporal order of 
message exchanges between objects. Used to illustrate operational flows and 
failure scenarios.